% SPDX-License-Identifier: GPL-3.0-or-later
% Copyright (C) 2026 Alessandro Greco
%% Testo alternativo per l'accessibilita' (solo EPUB).
\altfig{Configurazione di un servizio onion pubblico versione 3. Tor usa HiddenServiceDir /var/lib/tor/mio\_servizio/ e HiddenServicePort 80 127.0.0.1:8080. Nella directory genera hostname, hs\_ed25519\_public\_key e hs\_ed25519\_secret\_key: proteggere la chiave segreta, perché chi la ottiene può impersonare il servizio. L'indirizzo v3 contiene 56 caratteri base32 ottenuti dalla chiave pubblica Ed25519, da un checksum e dal byte di versione; il nome [v3].onion nel disegno è solo un segnaposto. Le connessioni all'indirizzo onion sulla porta virtuale 80 passano attraverso Tor e vengono inoltrate al server web in ascolto esclusivamente su 127.0.0.1:8080. Il servizio illustrato è pubblico e non richiede file alice.auth; l'autorizzazione dei client è una funzione separata.}
% Un servizio onion v3 pubblico: configurazione, identità e inoltro locale
\begin{tikzpicture}[font=\sffamily\scriptsize,
  file/.style={draw,rounded corners=1pt,fill=torgrigio!8,minimum height=0.55cm,minimum width=3.4cm,font=\ttfamily\scriptsize},
  box/.style={draw,rounded corners=2pt,align=center,minimum height=0.8cm}]
  \node[box,fill=torviola!10,text width=4.6cm,align=left,font=\ttfamily\scriptsize] (conf) at (0,1.6)
    {HiddenServiceDir /var/lib/tor/mio\_servizio/\\HiddenServicePort 80 127.0.0.1:8080};
  \node[font=\sffamily\tiny,text=torgrigio,above=1pt of conf] {torrc};
  \node[file] (h) at (0,0.1) {hostname};
  \node[file] (pk) at (0,-0.6) {hs\_ed25519\_public\_key};
  \node[file] (sk) at (0,-1.3) {hs\_ed25519\_secret\_key};
  \node[font=\sffamily\tiny,text=torgrigio,align=center] at (0,-1.95) {proteggerne accesso e chiave privata};
  \draw[->] (conf) -- (h);
  \node[box,fill=torverde!12,minimum width=2.4cm] (srv) at (6.4,-0.6) {server web\\127.0.0.1:8080};
  \node[box,fill=torviola!15,minimum width=2.4cm] (onion) at (6.4,1.2) {\texttt{[v3].onion}\\porta 80};
  \draw[->,torviola] (onion) -- node[right,font=\sffamily\tiny,align=left] {inoltro\\locale} (srv);
  \node[font=\sffamily\tiny,text=torgrigio] at (6.4,2.2) {v3: 56 base32};
  \draw[dashed,torgrigio,->] (pk.east) to[bend left=12] (onion.west);
  \node[anchor=west,font=\sffamily\tiny,text=avviso,align=left] at (4.4,-1.9) {chi ottiene la chiave segreta\\può impersonare il servizio};
  \draw[avviso,->] (sk.east) to[bend right=10] (4.3,-1.9);
\end{tikzpicture}
